\lstset{language=Python}
\lstset{frame=tb,
  language=Python,
  aboveskip=3mm,
  belowskip=3mm,
  showstringspaces=false,
  columns=flexible,
  basicstyle={\small\ttfamily},
  numbers=none,
  numberstyle=\tiny\color{gray},
  keywordstyle=\color{myteal},
  commentstyle=\color{myblue},
  stringstyle=\color{myyellow},
  breaklines=true,
  breakatwhitespace=true,
  tabsize=3
}

\chapter{Metodologija}
\label{chMe}

Okviren potek zagozdenja sistema:
\begin{enumerate}
\item{Generiranje dvdimenzionalne mreže $N$ točk z~Mitchellovim algoritmom~\cite{mitch}}.
\item{Postavitev $N$ elips z~ekscentričnostjo $e$ in naključnimi začetnimi orientacijami.}
\item{Implementacija prekrivalne funkcije~\cite{perram1985} za zaznavanje trkov med elipsami.}
\item{Postopno večanje elips in relaksacija vrtenja (Monte Carlo).}
\item{Analiza konfiguracij.}
\end{enumerate}
Dvodimenzionalnemu prostoru dodamo še periodične robne pogoje, ki jih vključimo na naslednji 
način:
\begin{lstlisting}[frame=single]  
    dx = point2 - point1
    if dx > 0.5*width:
        dx = dx - width
    elif dx < -0.5*width:
        dx = dx + width
\end{lstlisting}

\section{Mitchellov algoritem}

\begin{figure}[h!]
    \centering
    % GNUPLOT: LaTeX picture with Postscript
\begingroup
  \makeatletter
  \providecommand\color[2][]{%
    \GenericError{(gnuplot) \space\space\space\@spaces}{%
      Package color not loaded in conjunction with
      terminal option `colourtext'%
    }{See the gnuplot documentation for explanation.%
    }{Either use 'blacktext' in gnuplot or load the package
      color.sty in LaTeX.}%
    \renewcommand\color[2][]{}%
  }%
  \providecommand\includegraphics[2][]{%
    \GenericError{(gnuplot) \space\space\space\@spaces}{%
      Package graphicx or graphics not loaded%
    }{See the gnuplot documentation for explanation.%
    }{The gnuplot epslatex terminal needs graphicx.sty or graphics.sty.}%
    \renewcommand\includegraphics[2][]{}%
  }%
  \providecommand\rotatebox[2]{#2}%
  \@ifundefined{ifGPcolor}{%
    \newif\ifGPcolor
    \GPcolortrue
  }{}%
  \@ifundefined{ifGPblacktext}{%
    \newif\ifGPblacktext
    \GPblacktexttrue
  }{}%
  % define a \g@addto@macro without @ in the name:
  \let\gplgaddtomacro\g@addto@macro
  % define empty templates for all commands taking text:
  \gdef\gplbacktext{}%
  \gdef\gplfronttext{}%
  \makeatother
  \ifGPblacktext
    % no textcolor at all
    \def\colorrgb#1{}%
    \def\colorgray#1{}%
  \else
    % gray or color?
    \ifGPcolor
      \def\colorrgb#1{\color[rgb]{#1}}%
      \def\colorgray#1{\color[gray]{#1}}%
      \expandafter\def\csname LTw\endcsname{\color{white}}%
      \expandafter\def\csname LTb\endcsname{\color{black}}%
      \expandafter\def\csname LTa\endcsname{\color{black}}%
      \expandafter\def\csname LT0\endcsname{\color[rgb]{1,0,0}}%
      \expandafter\def\csname LT1\endcsname{\color[rgb]{0,1,0}}%
      \expandafter\def\csname LT2\endcsname{\color[rgb]{0,0,1}}%
      \expandafter\def\csname LT3\endcsname{\color[rgb]{1,0,1}}%
      \expandafter\def\csname LT4\endcsname{\color[rgb]{0,1,1}}%
      \expandafter\def\csname LT5\endcsname{\color[rgb]{1,1,0}}%
      \expandafter\def\csname LT6\endcsname{\color[rgb]{0,0,0}}%
      \expandafter\def\csname LT7\endcsname{\color[rgb]{1,0.3,0}}%
      \expandafter\def\csname LT8\endcsname{\color[rgb]{0.5,0.5,0.5}}%
    \else
      % gray
      \def\colorrgb#1{\color{black}}%
      \def\colorgray#1{\color[gray]{#1}}%
      \expandafter\def\csname LTw\endcsname{\color{white}}%
      \expandafter\def\csname LTb\endcsname{\color{black}}%
      \expandafter\def\csname LTa\endcsname{\color{black}}%
      \expandafter\def\csname LT0\endcsname{\color{black}}%
      \expandafter\def\csname LT1\endcsname{\color{black}}%
      \expandafter\def\csname LT2\endcsname{\color{black}}%
      \expandafter\def\csname LT3\endcsname{\color{black}}%
      \expandafter\def\csname LT4\endcsname{\color{black}}%
      \expandafter\def\csname LT5\endcsname{\color{black}}%
      \expandafter\def\csname LT6\endcsname{\color{black}}%
      \expandafter\def\csname LT7\endcsname{\color{black}}%
      \expandafter\def\csname LT8\endcsname{\color{black}}%
    \fi
  \fi
    \setlength{\unitlength}{0.0500bp}%
    \ifx\gptboxheight\undefined%
      \newlength{\gptboxheight}%
      \newlength{\gptboxwidth}%
      \newsavebox{\gptboxtext}%
    \fi%
    \setlength{\fboxrule}{0.5pt}%
    \setlength{\fboxsep}{1pt}%
    \definecolor{tbcol}{rgb}{1,1,1}%
\begin{picture}(7200.00,5040.00)%
    \gplgaddtomacro\gplbacktext{%
    }%
    \gplgaddtomacro\gplfronttext{%
    }%
    \gplbacktext
    \put(0,0){\includegraphics[width={360.00bp},height={252.00bp}]{./gnuplot/porazdelitev_1024}}%
    \gplfronttext
  \end{picture}%
\endgroup

    \caption[Blue noise]{Porazdelitev $1024$ točk.}
    \label{slika1}
\end{figure}

Okviren potek algoritma:
\begin{enumerate}
    \item{za začetni točki zgeneriramo naključni poziciji,}
    \item{zgeneriramo naključne pozicije kandidatov za naslednjo točko,}
    \item{izberemo tistega kandidata, ki je najdlje od vseh točk porazdelitve 
    (ima največjo minimalno oddaljenost),}
    \item{ponovimo korak 2 dokler ne dobimo željeno število točk.}
\end{enumerate}
Med algoritmom število kandidatov povečujemo sorazmerno s~številom že obstoječih
točk $N$. Pri nalogi smo tako na vsakem koraku generirali 
$ \left \lfloor{N/2}\right \rfloor + 1$ kandidatov.\\
%% napisi se nek del kako je to kul in primerjavo z modrim sumom blabla

\section{Eliptična kontaktna funkcija}

Kontaktna funkcija je dober kriterij za prekrivanje med dvema elipsama.\footnote{V~nalogi
uporabljamo dvodimenzionalni ekvivalent tridimenzionalne funkcije povzete 
po~\cite{perram1985}.}

\subsection{Ena elipsa}
Elipso $A$ definiramo s~funkcijo $E_a$, ki zadošča naslednjemu pogoju
\begin{equation}
    E_a (\vec{r}-\vec{r}_a, \theta_a) 
    \begin{cases}
        \quad <1 & \text{znotraj A}\\
        \quad =1 & \text{na površini A}\\
        \quad >1 & \text{zunaj A}
      \end{cases}. 
    \label{eq:cases}
\end{equation}
Izberemo

\begin{equation}
    E_a = (\vec{r} - \vec{r}_a)^{\top} \mathbf{A}^{-1} (\vec{r} - \vec{r}_a),
\label{eq:elipsa}
\end{equation}
kjer je $\vec{r}_a$ center elipse. Matriko $\mathbf{A}$ lahko zapišemo kot
\begin{equation}
    \mathbf{A} (\theta_a) = \mathbf{R}(\theta_a) \hat{\mathbf{A}} 
                            \mathbf{R}^{\top}(\theta_a),
\end{equation}
kjer je $R(\theta)$ rotacijska matrika, $\hat{\mathbf{A}}$ pa je definirana z~velikostjo 
polosi elipse $a_1$ in $a_2$ ter enotskih vektorjev $\vec{e}_i$ 
\begin{equation}
    \hat{\mathbf{A}} = \sum_{i=1}^{2} a_i^2 \vec{e}_i \vec{e}_i^{\top}.
\end{equation}
Matriki $\mathbf{A}$ in $\mathbf{A}^{-1}$ sta simetrični in pozitivno definitni.

\subsection{Prekrivanje dveh elips}
Definiramo funkcijo
\begin{equation}
    F(\vec{r}, \lambda) = \lambda E_a (\vec{r}) + (1 - \lambda) E_a (\vec{r}),
    \label{eq:f}
\end{equation}
ki je odvisna od pozicij $\vec{r}_a, \, \vec{r}_b$ in orientacij $\theta_a, \, \theta_b$
dveh elips $A$ ter $B$. Parameter $\lambda$ omejimo na interval $[0,1]$, tako, da je 
$F(\vec{r}, \lambda) \geq 0$. Pri fiksni vrednosti $\lambda$ ima $F(\vec{r}, \lambda)$
enolični minimum. Pri $\lambda=0$ je minimum $F=0$ pri $\vec{r} = \vec{r}_b$, pri 
$\lambda=1$ pa je minimum $F=0$ pri $\vec{r} = \vec{r}_a$. Za vse vmesne vrednosti
$\lambda$, je vrednost $\vec{r}$ pri minimumu $F(\vec{r}, \lambda)$ določena z
\begin{equation}
    \nabla F(\vec{r}, \lambda) = 0,
\end{equation}
oziroma
\begin{equation}
    \lambda \mathbf{A}^{-1} (\vec{r} - \vec{r}_a) + (1-\lambda) \mathbf{B}^{-1}
    (\vec{r} - \vec{r}_b) = 0.
\end{equation}
To lahko napišemo tudi kot
\begin{align}
    \vec{r}(\lambda) &- \vec{r}_a = (1-\lambda) \mathbf{A} \mathbf{C}^{-1} 
        \vec{r}_{a\,b}, \nonumber \\
    \vec{r}(\lambda) &- \vec{r}_b = - \lambda \mathbf{B} \mathbf{C}^{-1} \vec{r}_{a\,b}, 
    \label{eq:rab}
\end{align}
kjer je $\vec{r}_{a\,b} = \vec{r}_b - \vec{r}_a$ in $\mathbf{C}$ matrika
\begin{equation}
    \mathbf{C} = (1-\lambda) \mathbf{A} + \lambda \mathbf{B}.
    \label{eq:C}
\end{equation}
Rešitev~\ref{eq:rab} ustavimo v~\ref{eq:f}, pri čemer upoštevamo tudi~\ref{eq:C}
in definiramo funkcijo $f$ kot
\begin{equation}
    f(\lambda) = F (\vec{r}(\lambda), \lambda) = \lambda (1-\lambda)
        \vec{r}_{a\,b}^{\top} C^{-1} \vec{r}_{a\,b}.
\end{equation}

